\chapter{Évaluation et résultats}
\label{chap:evaluation}

\section*{Introduction}
\addcontentsline{toc}{section}{Introduction}

Ce chapitre présente le protocole d'évaluation et les résultats obtenus. Conformément à l'objectif
de fiabilité posé au chapitre~\ref{chap:contexte}, l'analyse des limites y occupe autant de place
que celle des performances : un chiffre favorable dont on ne connaît pas la portée n'apprend rien.

\section{Protocole}

\subsection{Jeu de référence}

Un jeu de référence de \textbf{\NbQuestions{} questions} a été construit à la main :
\textbf{\NbInScope{} dans le périmètre}, chacune associée au ou aux articles attendus et vérifiée
contre le texte réel du corpus, et \textbf{\NbOutScope{} hors périmètre} destinées à tester
l'abstention.

Le jeu a été construit en deux temps. Une première version de 37 questions a servi à l'essentiel
du développement. Un audit de couverture, mené une fois le système stabilisé, a révélé que
plusieurs comportements du code n'étaient exercés par \emph{aucune} question :

\begin{itemize}
  \item aucune question ne \textbf{nommait un article par son numéro}, alors que le module de
        recherche comporte un chemin dédié à ce cas (section~\ref{sec:refexplicite}), doublé d'un
        contournement du garde-fou. Cette branche n'était donc jamais testée ;
  \item les \textbf{articles~32 et~256}, seuls textes réintroduits à la main dans le corpus,
        n'apparaissaient dans aucune attente ;
  \item \textbf{trois Livres sur huit} --- organes de contrôle, conflits collectifs, dispositions
        finales, soit 60~articles --- n'étaient couverts par aucune question ;
  \item une seule question sur 32 attendait \textbf{plusieurs articles}, alors que c'est le cas
        courant en droit.
\end{itemize}

Vingt-cinq questions ont été ajoutées pour combler exactement ces manques
(tableau~\ref{tab:familles}). Le biais que cela introduit doit être énoncé : ces questions ciblent
des zones dont on \emph{savait} qu'elles existaient. Elles corrigent un angle mort, elles ne
rendent pas le jeu représentatif de l'usage réel.

\begin{table}[htbp]
  \centering
  \caption{Composition du jeu de référence par famille}
  \label{tab:familles}
  \small
  \begin{tabularx}{\textwidth}{@{}>{\raggedright\arraybackslash}p{3.6cm} c R@{}}
    \toprule
    \sffamily\bfseries\color{navy}Famille &
    \sffamily\bfseries\color{navy}Questions &
    \sffamily\bfseries\color{navy}Ce qu'elle exerce \\
    \midrule
    Périmètre d'origine     & 32 & Contrat, durée du travail, congés, maternité, délégués, sanctions, retraite \\
    Référence explicite     & 6  & Le chemin de recherche par numéro d'article \\
    Livre V                 & 4  & Inspection du travail, procès-verbaux, sanctions \\
    Livre VI                & 5  & Conflits collectifs, conciliation, arbitrage \\
    Livre VII               & 1  & Entrée en vigueur de la loi \\
    Articles multiples      & 3  & Réponses fondées sur plusieurs textes \\
    Articles rétablis       & 2  & Les articles 32 et 256, réintroduits à la main \\
    Hors périmètre          & 9  & Autres domaines juridiques, administratif, fiscal, hors droit \\
    \midrule
    \textbf{Total}          & \textbf{\NbQuestions} & \\
    \bottomrule
  \end{tabularx}
\end{table}

Ce jeu reste petit, et ce chapitre s'abstient donc de toute précision décimale sur les taux : à
cette échelle, un écart d'un point n'est pas un résultat, c'est du bruit.

\subsection{Scripts d'évaluation}

Trois scripts distincts, tous rejouables (tableau~\ref{tab:scripts}), produisent l'ensemble des
chiffres présentés ici. \textbf{Aucune valeur de ce chapitre n'a été saisie à la main} : les
nombres du texte comme les coordonnées des figures sont générés depuis les fichiers de mesure par
un script dédié (annexe~\ref{ann:depot}).

\begin{table}[htbp]
  \centering
  \caption{Scripts d'évaluation}
  \label{tab:scripts}
  \small
  % Première colonne à 3,7 cm : « run_answer_eval.py » est le plus long libellé
  % du rapport et débordait sur la colonne voisine à 2,7 cm.
  \begin{tabularx}{\textwidth}{@{}>{\raggedright\arraybackslash}p{3.7cm} R R@{}}
    \toprule
    \sffamily\bfseries\color{navy}Script &
    \sffamily\bfseries\color{navy}Mesure &
    \sffamily\bfseries\color{navy}Particularité \\
    \midrule
    \module{run\_eval.py}        & Recall@5, MRR
                                 & N'appelle pas le modèle : mesure la recherche seule \\
    \module{run\_answer\_eval.py} & Critères de succès
                                 & Citation, abstention, contrôle des citations \\
    \module{run\_measures.py}    & Latence, taxonomie
                                 & Produit également les données des figures \\
    \bottomrule
  \end{tabularx}
\end{table}

\subsection{Identité du run et reproductibilité}

Une précision méthodologique s'impose. La génération n'est pas déterministe : à température~0,1 le
modèle reste légèrement variable, et deux exécutions du même banc sur les mêmes questions ne
produisent pas exactement les mêmes réponses. Deux runs conservés dans le dépôt en témoignent :
l'un donne un refus à tort, l'autre trois.

\textbf{Tous les chiffres de ce chapitre proviennent d'un seul run}, celui du \DateMesure{} avec
le modèle \ModeleLLM{}, enregistré dans \module{eval/mesures-performance.json}. Aucun résultat
n'est composé à partir de plusieurs exécutions. Cette contrainte est plus sévère que nécessaire
pour un rapport de stage, mais elle évite le mélange silencieux de mesures issues de conditions
différentes --- lequel produit des synthèses flatteuses et fausses.

\begin{limite}
La conséquence à retenir est que les taux rapportés ici ne sont pas des constantes du système,
mais les valeurs d'une exécution. Établir des intervalles demanderait de rejouer le banc un grand
nombre de fois, ce qui n'a pas été fait : c'est une limite du protocole, pas une propriété des
chiffres.
\end{limite}

\section{Qualité de la recherche}

Les trois configurations ont été comparées sur les \NbInScope{} questions du périmètre, sans que
le modèle n'intervienne (figure~\ref{fig:recherche}).

\begin{figure}[htbp]
  \centering
  % Qualité de la recherche : les deux métriques sont bornées entre 0 et 1 et
% partagent donc un axe unique. Aucun axe secondaire.
%
% Barres horizontales, comme les autres graphiques du chapitre : l'étiquette
% de valeur se place au bout de chaque barre, ce qui rend toute collision
% impossible — en barres verticales, « 0,969 » et « 0,891 » se touchaient.
% Recall@5 et MRR par méthode de recherche
% Généré par rapport/build_data.py — NE PAS ÉDITER.
\newcommand{\RechercheRecall}{(0.849,0) (0.660,1) (0.868,2)}
\newcommand{\RechercheMrr}{(0.778,0) (0.535,1) (0.775,2)}

\begin{tikzpicture}
\begin{axis}[
  socle,
  xbar, bar width=8pt,
  height=5.2cm, width=0.72\linewidth,
  xmin=0, xmax=1.18,
  xtick={0,0.25,0.5,0.75,1},
  xmajorgrids, ymajorgrids=false,
  axis x line*=bottom, axis y line*=left,
  y axis line style={draw=none}, ytick style={draw=none},
  ytick={0,1,2}, y dir=reverse,
  yticklabels={{Dense seul},{BM25 seul},{Hybride (RRF)}},
  yticklabel style={font=\sffamily\small, text=ink, align=right},
  enlarge y limits=0.28,
  xlabel={valeur de la métrique},
  nodes near coords,
  % « fixed zerofill » : sans lui, 0,660 s'affiche « 0,66 » alors que le
  % tableau voisin écrit trois décimales.
  nodes near coords style={font=\sffamily\scriptsize, text=ink,
                           anchor=west, xshift=2.5pt,
                           /pgf/number format/fixed,
                           /pgf/number format/fixed zerofill,
                           /pgf/number format/precision=3},
  legend style={at={(0.5,1.03)}, anchor=south},
  legend image post style={line width=2.5pt},
]
% L'axe est inversé (y dir=reverse) : un décalage positif place la barre plus
% haut. Recall@5 vient donc en premier, dans l'ordre de la légende.
\addplot[draw=inptblue,  fill=inptblue,  fill opacity=0.9,  bar shift=5pt]
  coordinates {\RechercheRecall};
\addplot[draw=amberdark, fill=amberdark, fill opacity=0.9,  bar shift=-5pt]
  coordinates {\RechercheMrr};
\legend{Recall@5, MRR}
\end{axis}
\end{tikzpicture}

  \caption[Comparaison des méthodes de recherche]{Comparaison des trois méthodes de recherche sur
  les \NbInScope{} questions du périmètre. Recall@5 mesure la présence de l'article attendu dans
  les cinq premiers résultats, MRR sa position moyenne.}
  \label{fig:recherche}
\end{figure}

\subsection{Un résultat qui s'est inversé quand le jeu de test a grandi}
\label{sec:inversion}

Sur les \NbInScope{} questions du périmètre, \textbf{la recherche hybride dépasse la recherche
dense} en Recall@5 : \RecallHybride{} contre \RecallDense{}. Les deux méthodes sont en revanche
départagées à l'inverse sur le MRR --- \MrrDense{} pour le dense contre \MrrHybride{} pour
l'hybride --- soit un écart de trois millièmes, sans signification à cette taille d'échantillon.

Cette conclusion est l'exact contraire de celle que produisait la première version du jeu de
référence. Sur les 37 questions d'origine, les mêmes scripts donnaient un Recall@5 de 0,969 pour le
dense contre 0,875 pour l'hybride --- un écart net, en faveur du dense, qui avait été présenté
comme un résultat contre-intuitif du projet. Il n'en était pas un.

\begin{table}[htbp]
  \centering
  \caption{Effet de l'élargissement du jeu de référence sur la comparaison des méthodes}
  \label{tab:inversion}
  \small
  \begin{tabularx}{\textwidth}{@{}>{\raggedright\arraybackslash}p{3.4cm} c c R@{}}
    \toprule
    \sffamily\bfseries\color{navy}Méthode &
    \sffamily\bfseries\color{navy}37 questions &
    \sffamily\bfseries\color{navy}\NbQuestions{} questions &
    \sffamily\bfseries\color{navy}Lecture \\
    \midrule
    Dense seul     & \textbf{0,969} & \RecallDense{} & Chute lourdement \\
    BM25 seul      & 0,781 & \RecallBmXXV{} & Chute aussi \\
    Hybride (RRF)  & 0,875 & \textbf{\RecallHybride{}} & \textbf{Seule méthode à peu près stable} \\
    \bottomrule
  \end{tabularx}
\end{table}

L'explication tient entièrement à la \emph{composition} du jeu d'origine. Ses 32 questions étaient
toutes des reformulations en langage courant, sans terme technique repris du texte de loi : la
configuration exactement favorable aux embeddings et défavorable à BM25. Le dense y gagnait donc,
non parce qu'il est meilleur, mais parce qu'on ne lui posait que les questions qu'il sait traiter.

Les vingt-cinq questions ajoutées introduisent ce que le jeu d'origine ne contenait pas : des
références explicites à des numéros d'articles, et le vocabulaire figé des Livres~V et~VI ---
«~procès-verbal~», «~mise en demeure~», «~commission provinciale d'enquête et de conciliation~».
Ce sont précisément les correspondances littérales sur lesquelles le chapitre~\ref{chap:fondements}
annonçait que le lexical l'emporterait. La prédiction théorique était juste ; c'est la mesure qui
était biaisée.

\begin{constat}
Le résultat à retenir n'est pas «~l'hybride gagne~» --- l'écart reste faible, et un troisième jeu
de questions le déplacerait encore. C'est que \textbf{la première conclusion mesurait le jeu de
test, pas le système}. Un banc d'évaluation ne devient un instrument de mesure qu'une fois vérifié
qu'il exerce réellement toutes les capacités du système ; jusque-là, il mesure surtout les
hypothèses de celui qui l'a écrit.

\smallskip
La conséquence pratique est que l'hybride est conservé. Il est la seule des trois configurations à
ne pas s'effondrer quand la nature des questions change, ce qui est la propriété qui compte pour un
système destiné à des questions qu'on ne choisit pas.
\end{constat}

\section{Qualité des réponses}

\subsection{Taxonomie des résultats}

Chaque réponse a été classée dans une catégorie unique (figure~\ref{fig:taxonomie}).

\begin{figure}[htbp]
  \centering
  % Taxonomie des réponses. Les catégories sont des états, pas des séries :
% la couleur suit donc une échelle de statut, du vert (résultat visé) au
% rouge (échec), et chaque barre porte son libellé — l'identité ne repose
% jamais sur la seule couleur.
%
% Deux réglages non évidents :
%   - les libellés d'axe sont écrits en clair : pgfplots ne développe pas une
%     macro passée à « yticklabels », quelle que soit la manière de la lui
%     donner. Leur ORDRE doit rester celui de TAXONOMY_ORDER dans
%     rapport/build_data.py, qui fixe l'ordre des barres ;
%   - « bar shift=0 » : en mode xbar, chaque \addplot reçoit sinon un décalage
%     propre et les barres se retrouvent échelonnées au lieu d'être alignées
%     sur leur ordonnée.
% Taxonomie des réponses générées, du meilleur au pire résultat
% Généré par rapport/build_data.py — NE PAS ÉDITER.
\newcommand{\TaxonomieLabels}{{Sourcée, article attendu},{Sourcée, autre article},{Abstention correcte},{Citation non vérifiée},{Sans citation},{Refus à tort}}
\newcommand{\TaxonomieBarA}{(41,0)}
\newcommand{\TaxonomieBarB}{(2,1)}
\newcommand{\TaxonomieBarC}{(9,2)}
\newcommand{\TaxonomieBarD}{(6,3)}
\newcommand{\TaxonomieBarE}{(2,4)}
\newcommand{\TaxonomieBarF}{(2,5)}
\newcommand{\TaxonomieMax}{41}
\newcommand{\TaxonomieTotal}{62}

\begin{tikzpicture}
\begin{axis}[
  barresH,
  height=6cm,
  % La largeur de pgfplots ne compte que l'aire tracée : les libellés de
  % catégorie s'y ajoutent. On la réduit pour que l'ensemble tienne dans la
  % justification.
  width=0.74\linewidth,
  xmin=0, xmax=\TaxonomieMax*1.16,
  ytick={0,...,5},
  yticklabels={{Sourcée, article attendu},{Sourcée, autre article},
               {Abstention correcte},{Citation non vérifiée},
               {Sans citation},{Refus à tort}},
  yticklabel style={font=\sffamily\scriptsize, text=ink, align=right},
  xlabel={nombre de questions \textcolor{slate}{(sur \TaxonomieTotal)}},
  nodes near coords,
  bar width=11pt,
  every axis plot/.append style={bar shift=0},
]
\addplot[draw=leafgreen, fill=leafgreen, fill opacity=0.88] coordinates {\TaxonomieBarA};
\addplot[draw=leafgreen, fill=leafgreen, fill opacity=0.45] coordinates {\TaxonomieBarB};
\addplot[draw=inptblue,  fill=inptblue,  fill opacity=0.85] coordinates {\TaxonomieBarC};
\addplot[draw=amberdark, fill=amberdark, fill opacity=0.85] coordinates {\TaxonomieBarD};
\addplot[draw=amberdark, fill=amberdark, fill opacity=0.5]  coordinates {\TaxonomieBarE};
\addplot[draw=anrtred,   fill=anrtred,   fill opacity=0.85] coordinates {\TaxonomieBarF};
\end{axis}

% La catégorie absente est le vrai résultat : aucune réponse produite hors
% corpus. Elle est posée sous la figure, hors de l'axe, pour ne pas se
% confondre avec une barre de valeur nulle.
\node[font=\sffamily\scriptsize, text=leafgreen!40!black, anchor=north west,
      align=left, inner sep=0pt]
  at ($(current bounding box.south west)+(0,-3mm)$)
  {\textbf{Répond alors qu'il devrait s'abstenir : 0} \textcolor{slate}{— catégorie vide}};
\end{tikzpicture}

  \caption[Classement des réponses générées]{Classement des \NbQuestions{} réponses générées. La
  catégorie la plus importante est celle qui n'apparaît pas : aucune réponse n'a été produite sur
  un sujet absent du corpus.}
  \label{fig:taxonomie}
\end{figure}

Pour un outil juridique, c'est bien l'absence de la dernière catégorie qui compte. Un système qui
répondrait à côté sur une question qu'il ne couvre pas serait dangereux d'une manière que ne
compenserait aucun taux de réussite par ailleurs.

\subsection{Indicateurs de fiabilité}

Quatre indicateurs résument la fiabilité des réponses (figure~\ref{fig:fiabilite}). Ils ne portent
pas sur les mêmes populations, ce que la figure rend explicite.

\begin{figure}[htbp]
  \centering
  % Indicateurs de fiabilité. Une seule série, donc une seule teinte et pas de
% légende : le titre de l'axe suffit à la nommer. Le rail clair rappelle que
% l'échelle va jusqu'à 100 %, ce qu'une barre nue laisse deviner.
% Indicateurs de fiabilité, en pourcentage
% Généré par rapport/build_data.py — NE PAS ÉDITER.
\newcommand{\FiabiliteLabels}{{Citation systématique},{Citations vérifiées},{Article attendu cité},{Abstention hors périmètre}}
\newcommand{\FiabiliteCoords}{(96.1,0) (81.1,1) (88.2,2) (100.0,3)}
\newcommand{\FiabiliteDetailA}{49/51}
\newcommand{\FiabiliteDetailB}{77/95}
\newcommand{\FiabiliteDetailC}{45/51}
\newcommand{\FiabiliteDetailD}{9/9}

\begin{tikzpicture}[
  detail/.style={font=\sffamily\scriptsize, text=slate, anchor=west}]
\begin{axis}[
  barresH,
  height=4.8cm,
  % Largeur réduite : les libellés à gauche et les effectifs à droite
  % s'ajoutent à l'aire tracée.
  width=0.62\linewidth,
  xmin=0, xmax=134,
  xtick={0,25,50,75,100},
  xticklabel={\pgfmathprintnumber{\tick}\,\%},
  ytick={0,...,3},
  % Libellés en clair : pgfplots ne développe pas une macro donnée à
  % « yticklabels ». L'ordre suit celui de la liste « fiab » dans
  % rapport/build_data.py, qui fixe l'ordre des barres.
  yticklabels={{Citation systématique},{Citations vérifiées},
               {Article attendu cité},{Abstention hors périmètre}},
  yticklabel style={font=\sffamily\scriptsize, text=ink, align=right},
  xlabel={part des cas conformes},
  bar width=11pt,
  every axis plot/.append style={bar shift=0},
]
% Rail de fond : la portion non atteinte.
\addplot[draw=none, fill=mist, forget plot]
  coordinates {(100,0) (100,1) (100,2) (100,3)};
\addplot[draw=inptblue, fill=inptblue, fill opacity=0.9,
         nodes near coords={\pgfmathprintnumber[precision=0]{\pgfplotspointmeta}\,\%},
         nodes near coords style={font=\sffamily\scriptsize\bfseries, text=ink,
                                  anchor=west, xshift=3pt}]
  coordinates {\FiabiliteCoords};

% Effectif réel derrière chaque pourcentage : 93 % sur 45 citations n'a pas
% le même poids que 100 % sur 5 questions, et la figure doit le dire.
\node[detail] at (axis cs:116,0) {\FiabiliteDetailA};
\node[detail] at (axis cs:116,1) {\FiabiliteDetailB};
\node[detail] at (axis cs:116,2) {\FiabiliteDetailC};
\node[detail] at (axis cs:116,3) {\FiabiliteDetailD};
\end{axis}
\end{tikzpicture}

  \caption[Indicateurs de fiabilité]{Indicateurs de fiabilité, avec l'effectif réel derrière
  chaque taux. Les dénominateurs diffèrent : \TauxAbstention~\% sur \NbOutScope{} questions et
  \TauxVerif~\% sur \NbCitations{} citations n'ont pas le même poids de preuve.}
  \label{fig:fiabilite}
\end{figure}

Sur les \NbRepondues{} questions du périmètre auxquelles le système a effectivement répondu,
\TauxCitation~\% citent au moins un article et \TauxSourceAttendue~\% citent l'article attendu.
Sur les \NbCitations{} citations produites au total, \NbNonVerifiees{} ne correspondent à aucun
article présent dans le contexte fourni, soit \TauxVerif~\% de citations vérifiées. Les
\NbOutScope{} questions hors périmètre ont toutes été refusées.

Ce taux de \TauxVerif~\% appelle un examen, car il est nettement inférieur à celui qu'affichait le
jeu d'origine (93~\% sur 45~citations). Le nombre de citations produites a plus que doublé
(\NbCitations{} contre 45), et le nombre de citations signalées a été multiplié par six.

Le mécanisme principal reste le \textbf{renvoi interne} : le texte d'un article fourni mentionne
lui-même d'autres articles («~\dots dans les conditions prévues par les articles~154 et~156~»),
que le modèle reprend. Le contrôle les signale parce qu'ils ne figurent pas dans le contexte ; il
sur-signale volontairement, plutôt que de risquer de manquer une véritable invention.

\begin{limite}
\textbf{Cette explication ne couvre cependant qu'une partie des cas, et le banc ne permet pas de
statuer sur le reste.} Un recoupement automatique des \NbNonVerifiees{} citations signalées avec
le texte des articles fournis confirme le renvoi interne dans six cas seulement. Les autres restent
indéterminés --- notamment une réponse créditée des articles~1 à~6, une suite trop régulière pour
être une citation juridique, et qui évoque plutôt un faux positif du détecteur.

\smallskip
La cause de cette incertitude est une lacune du banc lui-même : \module{run\_measures.py}
enregistrait les numéros cités mais \emph{pas le texte de la réponse}, ce qui rend tout diagnostic
a posteriori impossible. Le script a été corrigé pour conserver les réponses, mais la présente
campagne de mesure a été exécutée avant cette correction. Le chiffre de \TauxVerif~\% doit donc
être lu comme un \textbf{plancher} : il compte comme suspecte toute citation non retrouvée, y
compris celles que le détecteur a lui-même mal extraites.
\end{limite}

\section{Le garde-fou d'abstention}

\subsection{Calibration du seuil}

Le seuil d'abstention a été calibré empiriquement, non choisi arbitrairement. La mesure des scores
de recherche sur les deux populations de questions fait apparaître trois bandes
(figure~\ref{fig:seuil}) :

\begin{itemize}
  \item questions de droit du travail : de \ScoreInMin{} à \ScoreInMax{} ;
  \item questions d'un \emph{autre domaine juridique ou administratif} : de
        \ScoreJuridiqueMin{} à \ScoreJuridiqueMax{} --- plus hautes que la bande suivante car
        elles partagent le vocabulaire juridique français ;
  \item questions sans rapport avec le droit : de \ScoreHorsDroitMin{} à \ScoreHorsDroitMax{}.
\end{itemize}

\begin{figure}[htbp]
  \centering
  % Score de recherche de chacune des 37 questions, réparti selon le périmètre.
% Le point de la figure est la position du seuil : il sépare proprement la
% bande basse, et pas les deux bandes hautes — c'est délibéré.
%
% Les trois bandes forment une série catégorielle (bleu, ambre, prune) et non
% une échelle de statut : une question hors périmètre n'est pas un échec, elle
% relève d'un autre domaine. Le vert et le rouge, réservés au succès et à
% l'échec ailleurs dans le rapport, sont donc écartés ici.
% Score de recherche par question, réparti en trois bandes
% Généré par rapport/build_data.py — NE PAS ÉDITER.
\newcommand{\ScoresTravail}{(0.5955,2.055) (0.6039,1.890) (0.6053,1.945) (0.6131,2.000) (0.6328,2.055) (0.6542,2.110) (0.6611,1.890) (0.6716,1.945) (0.6743,2.000) (0.6793,2.055) (0.6797,2.110) (0.6920,1.890) (0.6924,1.945) (0.6938,2.000) (0.6946,2.055) (0.6958,2.110) (0.6960,1.890) (0.7050,1.945) (0.7076,2.000) (0.7092,2.055) (0.7124,2.110) (0.7135,1.890) (0.7153,1.945) (0.7160,2.000) (0.7164,2.055) (0.7173,2.110) (0.7180,1.890) (0.7214,1.945) (0.7263,2.000) (0.7278,2.055) (0.7335,2.110) (0.7373,1.890) (0.7434,1.945) (0.7487,2.000) (0.7492,2.055) (0.7550,2.110) (0.7558,1.890) (0.7640,1.945) (0.7658,2.000) (0.7670,2.055) (0.7674,2.110) (0.7707,1.890) (0.7805,1.945) (0.7826,2.000) (0.7840,2.055) (0.7881,2.110) (0.7935,1.890)}
\newcommand{\ScoresJuridique}{(0.4453,1.055) (0.5147,0.890) (0.5150,0.945) (0.5561,1.000) (0.6027,1.110)}
\newcommand{\ScoresHorsDroit}{(0.3071,-0.110) (0.3357,-0.055) (0.4319,0.000) (0.4461,0.110)}

\begin{tikzpicture}
\begin{axis}[
  bandesY,
  height=5.8cm,
  xmin=0.26, xmax=0.85,
  ymin=-0.75, ymax=3.05,
  xtick={0.3,0.4,0.5,0.6,0.7,0.8},
  xticklabel style={/pgf/number format/fixed,
                    /pgf/number format/precision=1},
  xlabel={score de recherche \textcolor{slate}{(similarité cosinus maximale)}},
  ytick={0,1,2},
  yticklabels={},
  legend style={at={(0.5,1.1)}, anchor=south},
]
% Seuil : filet vertical plein hauteur, étiquette posée en haut pour ne
% croiser ni les graduations ni les libellés de zone.
\draw[anrtred, line width=0.8pt, dash pattern=on 3pt off 2pt]
  (axis cs:0.42,-0.75) -- (axis cs:0.42,3.05);
\node[font=\sffamily\scriptsize\bfseries, text=anrtred, anchor=south east]
  at (axis cs:0.415,3.05) {seuil \SeuilAbstention~};

\addplot[only marks, mark=*, mark size=1.9pt,
         draw=inptblue, fill=inptblue, fill opacity=0.7]
  coordinates {\ScoresTravail};
\addplot[only marks, mark=*, mark size=1.9pt,
         draw=amberdark, fill=amberdark, fill opacity=0.8]
  coordinates {\ScoresJuridique};
\addplot[only marks, mark=*, mark size=1.9pt,
         draw=plum, fill=plum, fill opacity=0.8]
  coordinates {\ScoresHorsDroit};
\legend{droit du travail, autre domaine juridique, hors du droit}

% Étiquettes de bande, à gauche du nuage : elles remplacent les graduations
% de l'axe vertical, qui ne porte aucune quantité.
\node[font=\sffamily\scriptsize, text=inptblue!70!black, anchor=east]
  at (axis cs:0.272,2) {\NbInScope{} questions};
\node[font=\sffamily\scriptsize, text=amberdark!70!black, anchor=east]
  at (axis cs:0.272,1) {3 questions};
\node[font=\sffamily\scriptsize, text=plum!70!black, anchor=east]
  at (axis cs:0.272,0) {2 questions};

% Les deux zones que le seuil sépare réellement, sous le nuage.
\node[font=\sffamily\scriptsize\itshape, text=anrtred, anchor=north east]
  at (axis cs:0.412,-0.42) {rejeté avant le modèle};
\node[font=\sffamily\scriptsize\itshape, text=slate, anchor=north west]
  at (axis cs:0.428,-0.42) {transmis au modèle, qui peut encore s'abstenir};
\end{axis}
\end{tikzpicture}

  \caption[Séparation des scores de recherche]{Score de recherche par question. Les six questions
  nommant un article par son numéro sont exclues : leur score est \emph{forcé} à 1,0 par le module
  de recherche (section~\ref{sec:refexplicite}) et ne mesure aucune similarité. Les bandes se
  chevauchent désormais, ce qui n'était pas visible sur le jeu d'origine.}
  \label{fig:seuil}
\end{figure}

\begin{limite}
\textbf{Les bandes se recouvrent, et aucun seuil ne peut donc les séparer.} La question
«~Quelles sont les conditions pour obtenir un passeport marocain~?~» obtient \ScoreJuridiqueMax{},
\emph{au-dessus} de la question de droit du travail la moins bien classée (\ScoreInMin{}, sur la
durée maximale d'une mission d'intérim). Il n'existe aucune valeur de seuil qui accepterait la
seconde tout en rejetant la première.

\smallskip
Le jeu d'origine, avec ses cinq seules questions hors périmètre, laissait croire à une séparation
nette. Elle était un artefact de l'échantillon --- le même biais que celui décrit en
section~\ref{sec:inversion}, sur une autre grandeur.
\end{limite}

Ce constat renforce, en le durcissant, l'argument de conception du
chapitre~\ref{chap:conception} : \textbf{le seuil ne doit pas chercher à trancher la pertinence
juridique}. Il n'en est pas capable, et le jour où le corpus s'étendra à un second code, ces
questions devront de toute façon devenir répondables.

Son rôle réel est plus modeste et parfaitement assumé : écarter à moindre coût la bande clairement
non pertinente, avant d'engager l'inférence. Sur \NbOutScope{} questions hors périmètre, il n'en
arrête que \NbGarde{} --- les sept autres franchissent le seuil et sont rattrapées en aval par
l'abstention au niveau du prompt, qui les refuse toutes. C'est bien la \emph{combinaison} des deux
garde-fous, et non le seuil seul, qui produit les \TauxAbstention~\% d'abstention correcte.

\subsection{Coût de l'abstention}

\NbRefusTort{} questions valides sur \NbInScope{} ont été refusées à tort. C'est le prix d'une
abstention stricte, et l'arbitrage est assumé : pour un outil juridique, un refus injustifié est
préférable à une réponse inventée. Le premier laisse l'usager où il était ; le second l'égare en
lui donnant confiance.

Cet équilibre reste néanmoins perfectible, et les deux cas concernés sont instructifs. Leurs scores
de recherche --- 0,672 et 0,695 --- sont largement au-dessus du seuil : \emph{aucun des deux n'a
été arrêté par le garde-fou}. Le refus vient donc du \textbf{modèle}, non de la recherche,
c'est-à-dire de l'effet des règles ajoutées en section~\ref{sec:hallucination}. La marge de progrès
est du côté du prompt, pas du seuil.

Le second cas est le plus révélateur : à la question «~Comment l'arbitre est-il désigné en cas de
conflit collectif~?~», l'article~568 attendu \emph{figurait} dans le contexte fourni, et le modèle
a néanmoins refusé sans citer aucun article. Le sujet relève du Livre~VI, absent du jeu de
référence d'origine : c'est une famille de questions sur laquelle le prompt n'avait jamais été
éprouvé.

\section{Latence}

Le temps de réponse a été mesuré étape par étape sur les \NbQuestions{} questions du jeu de
référence (figure~\ref{fig:latence}).

\begin{figure}[htbp]
  \centering
  % Latence, en deux lectures : de quoi le temps de réponse est fait, puis
% comment il se distribue d'une question à l'autre.
% Latence : composition médiane et distribution par question
% Généré par rapport/build_data.py — NE PAS ÉDITER.
\newcommand{\LatMedRecherche}{0.164}
\newcommand{\LatMedGeneration}{3.827}
\newcommand{\LatMedAbstention}{0.134}
\newcommand{\LatCourbeTotal}{(1,1.179) (2,1.678) (3,1.731) (4,1.846) (5,1.872) (6,1.898) (7,1.977) (8,2.07) (9,2.082) (10,2.093) (11,2.146) (12,2.249) (13,2.374) (14,2.377) (15,2.406) (16,2.437) (17,2.536) (18,2.652) (19,2.709) (20,2.795) (21,2.896) (22,2.987) (23,3.009) (24,3.209) (25,3.267) (26,3.388) (27,3.495) (28,3.552) (29,3.914) (30,3.919) (31,4.044) (32,4.117) (33,4.389) (34,4.495) (35,4.511) (36,4.584) (37,4.762) (38,4.826) (39,4.98) (40,5.084) (41,5.286) (42,5.295) (43,5.317) (44,5.37) (45,6.15) (46,6.329) (47,6.335) (48,6.455) (49,7.14) (50,8.408) (51,9.227) (52,10.054) (53,10.095) (54,10.257) (55,10.636) (56,10.784) (57,10.901) (58,11.521) (59,14.896) (60,16.315)}
\newcommand{\LatCourbeRecherche}{(1,0.147) (2,0.165) (3,0.171) (4,0.167) (5,0.151) (6,0.156) (7,0.154) (8,0.149) (9,0.16) (10,0.175) (11,0.154) (12,0.159) (13,0.141) (14,0.352) (15,0.156) (16,0.166) (17,0.179) (18,0.168) (19,0.16) (20,0.161) (21,0.184) (22,0.171) (23,0.246) (24,0.155) (25,0.145) (26,0.171) (27,0.176) (28,0.148) (29,0.166) (30,0.132) (31,0.177) (32,0.149) (33,0.177) (34,0.177) (35,0.141) (36,0.173) (37,0.145) (38,0.169) (39,0.194) (40,0.245) (41,0.16) (42,0.17) (43,0.164) (44,0.164) (45,0.151) (46,0.146) (47,0.182) (48,0.174) (49,0.14) (50,0.151) (51,0.181) (52,0.208) (53,0.238) (54,0.205) (55,0.148) (56,0.212) (57,0.147) (58,0.191) (59,0.153) (60,0.163)}
\newcommand{\LatNbCourbe}{60}

\begin{tikzpicture}

% ------------------------------------------- (a) composition du temps médian
\begin{axis}[
  name=compo,
  socle,
  xbar stacked,
  height=3.2cm, width=0.9\linewidth,
  bar width=12pt,
  xmin=0, xmax=4.3,
  xmajorgrids, axis y line=none,
  ytick={0,1}, y dir=reverse,
  yticklabels={},
  enlarge y limits=0.75,
  xlabel={secondes},
  title={\textbf{(a)}~~Composition du temps de réponse médian},
  title style={at={(0,1.32)}, anchor=south west},
  legend style={at={(1,1.32)}, anchor=south east},
]
\addplot[draw=white, line width=0.7pt, fill=inptblue, fill opacity=0.9]
  coordinates {(\LatMedRecherche,0) (\LatMedAbstention,1)};
\addplot[draw=white, line width=0.7pt, fill=amberdark, fill opacity=0.85]
  coordinates {(\LatMedGeneration,0) (0,1)};
\legend{recherche \LatRecherche{} s, génération \LatGeneration{} s}

\node[font=\sffamily\scriptsize\bfseries, text=ink, anchor=west]
  at (axis cs:3.52,0) {\LatTotale{} s};
\node[font=\sffamily\scriptsize\bfseries, text=anrtred, anchor=west]
  at (axis cs:0.24,1) {\LatAbstention{} s};

\node[font=\sffamily\scriptsize, text=slate, anchor=east]
  at (axis cs:-0.05,0) {réponse complète};
\node[font=\sffamily\scriptsize, text=slate, anchor=east]
  at (axis cs:-0.05,1) {abstention};
\end{axis}

% ------------------------------------- (b) distribution question par question
\begin{axis}[
  socle,
  at={($(compo.south west)-(0,2.5cm)$)}, anchor=north west,
  height=5.4cm, width=0.9\linewidth,
  xmin=0.5, xmax=\LatNbCourbe+0.5,
  ymin=0, ymax=15,
  ytick={0,5,10},
  ymajorgrids,
  xlabel={les \LatNbCourbe{} questions transmises au modèle, triées par temps total},
  ylabel={secondes},
  title={\textbf{(b)}~~Distribution du temps de réponse},
  title style={at={(0,1.06)}, anchor=south west},
]
\addplot[ybar, bar width=4.2pt, draw=none, fill=amberdark, fill opacity=0.6]
  coordinates {\LatCourbeTotal};
\addplot[ybar, bar width=4.2pt, draw=none, fill=inptblue]
  coordinates {\LatCourbeRecherche};

\draw[slate, line width=0.7pt, dash pattern=on 3pt off 2pt]
  (axis cs:0.5,3.486) -- (axis cs:\LatNbCourbe+0.5,3.486);
\node[font=\sffamily\scriptsize, text=slate, anchor=south west]
  at (axis cs:1,3.6) {médiane \LatTotale{} s};

\node[font=\sffamily\scriptsize, text=slate, anchor=north west, align=left,
      text width=52mm]
  at (axis cs:1.5,14.6) {la part bleue reste plate d'un bout à l'autre :
                         tout l'écart vient de l'inférence};
\draw[slate, line width=0.5pt, -{Stealth[length=1.8mm]}]
  (axis cs:11,11.2) to[out=-25,in=150] (axis cs:20,0.9);
\end{axis}

\end{tikzpicture}

  \caption[Répartition du temps de réponse]{Répartition du temps de réponse.
  \textbf{(a)}~composition du temps médian ; \textbf{(b)}~distribution sur les questions
  effectivement transmises au modèle. Les deux panneaux partagent le même code couleur :
  bleu pour la recherche, ambre pour la génération.}
  \label{fig:latence}
\end{figure}

Deux observations méritent d'être relevées.

\paragraph{La recherche ne représente que \PartRecherche~\% du temps total}
L'ensemble de la chaîne de recherche --- vectorisation de la question, interrogation des deux
index, fusion --- coûte \LatRecherche~seconde. Tout le reste est de l'inférence. Optimiser
davantage la recherche n'apporterait donc rien au temps de réponse perçu : c'est une conclusion
utile, car l'intuition d'ingénieur pousse spontanément à optimiser la partie du système qu'on a
écrite soi-même.

\paragraph{Le garde-fou divise le temps de réponse par \FacteurGarde{}}
Une question arrêtée avant le modèle reçoit sa réponse en \LatAbstention~seconde, contre
\LatTotale~secondes pour une réponse complète. Le placement du garde-fou \emph{avant} le modèle,
décidé pour des raisons de fiabilité au chapitre~\ref{chap:conception}, produit ainsi un bénéfice
de performance qui n'était pas son objectif.

La dispersion est par ailleurs importante : de 1,2 à \LatMax~secondes selon la longueur de la
réponse produite. La médiane est donc un indicateur plus honnête que la moyenne, que quelques
réponses longues suffisent à déplacer.

\section{Défaillance observée et corrigée}
\label{sec:hallucination}

L'évaluation a mis au jour une défaillance qui n'était pas apparue lors des tests manuels. Elle
est rapportée ici en détail, car elle constitue le mode d'échec le plus dangereux pour ce type de
système.

Interrogé sur la procédure de divorce --- question hors périmètre --- le modèle refusait
correctement, \emph{puis inventait une réponse complète} citant un «~Code civil marocain
(Loi~14-82)~» et son «~article~96~», en affirmant l'avoir consulté. Ces références n'existent
pas : le droit de la famille marocain relève de la Moudawana (loi~70-03), et aucun texte de ce
type n'avait été fourni au modèle.

\begin{limite}
Ce comportement --- reconnaître l'absence d'information, puis produire malgré tout un contenu
inventé --- est \textbf{plus trompeur qu'une invention franche}. La prudence apparente du début de
la réponse renforce la crédibilité de ce qui suit : le lecteur y voit un système qui connaît ses
limites, donc digne de confiance sur le reste.
\end{limite}

La correction a consisté à interdire explicitement, dans le prompt système, toute mention d'un
numéro de loi ou d'article provenant d'un autre code. Une première version de cette règle s'est
révélée \textbf{trop stricte} : le taux d'abstention correcte est passé à \TauxAbstention~\%, mais
les refus à tort ont triplé. Une règle de rééquilibrage --- \emph{par défaut, réponds} --- a permis
de conserver le bénéfice sans en payer tout le coût.

Cet épisode illustre l'intérêt d'un banc rejouable. Sans lui, la sur-correction serait passée
inaperçue : le système aurait simplement paru «~plus prudent~», alors qu'il devenait moins utile.
Une dégradation qui prend l'apparence d'une vertu est précisément le genre de régression qu'aucun
test manuel ne détecte.

\section{Limite de mesure assumée}

Le contrôle des citations vérifie qu'un article cité \textbf{figurait dans le contexte fourni} ---
pas qu'il \textbf{dise réellement} ce que la réponse lui attribue. Une réponse citant un article
réel tout en en déformant le contenu est comptée ici comme sourcée.

Détecter ce cas demanderait une vérification d'implication (\anglais{entailment}) ou une relecture
humaine systématique. C'est hors du périmètre de ce stage, et le chiffre de \TauxVerif~\% de
citations vérifiées doit être lu avec cette réserve : il mesure la \emph{traçabilité} des sources,
non la \emph{fidélité sémantique} des affirmations. C'est un plancher de confiance, pas une
garantie d'exactitude.

\section*{Conclusion}
\addcontentsline{toc}{section}{Conclusion}

Le système atteint les objectifs de fiabilité fixés : citation quasi systématique
(\TauxCitation~\%), abstention correcte sur l'intégralité des questions hors périmètre, et aucune
réponse produite sur un sujet absent du corpus.

Trois résultats vont à l'encontre des attentes initiales, et aucun n'était accessible par
l'intuition. L'hallucination inter-codes, d'abord, invisible aux tests manuels. Mais surtout, les
deux constats produits par l'élargissement du jeu de référence : \textbf{la comparaison des
méthodes de recherche s'est inversée} (section~\ref{sec:inversion}), et \textbf{les bandes de score
d'abstention se recouvrent}, alors que le jeu d'origine les montrait nettement séparées.

Ces deux derniers points ont une leçon commune, et c'est peut-être le principal acquis
méthodologique de ce stage : les deux conclusions renversées n'étaient pas fausses parce que le
système avait changé --- il n'avait pas changé --- mais parce que le \emph{jeu de test} ne
l'exerçait pas complètement. Un banc d'évaluation dont on n'a pas vérifié la couverture mesure les
hypothèses de son auteur autant que le système qu'il prétend juger.
